% §2.1 The Binary Symmetric Channel
%
% Cross-reference resolution status:
%
% EXPECTED UNRESOLVED (forward refs that will resolve in later Plans):
%   \Cref{sec:three-faces}       -- chapter 1 §1.1 (will resolve in Plan 4)
%   \Cref{sec:asymmetric-channel}-- chapter 2 §2.2 (will resolve in Task 3 of Plan 2)
%   \Cref{sec:channel-composition}-- chapter 2 §2.3 (will resolve in Task 4 of Plan 2)
%   \Cref{sec:model-order}       -- chapter 2 §2.4 (will resolve in Task 5 of Plan 2)
%   \ref{ex:bsc-capacity}        -- §2.1 exercise on BSC capacity (will resolve in Task 7)
%
% ALREADY RESOLVED (chapter 4, drafted in Plan 1):
%   \Cref{ch:axioms-kronecker}   -- chapter 4 label
%   \Cref{sec:kronecker}         -- chapter 4 §4.4
%   \Cref{sec:axiom-2}           -- chapter 4 §4.3
%
% Build warnings for the unresolved labels are expected at this stage.
% They will be resolved when Plan 2 Tasks 3-5 and Task 7 land and Plan 4 runs.

\section{The Binary Symmetric Channel}
\label{sec:bsc}

\subsection*{One bit, one channel}

\Cref{sec:three-faces} introduced three framings of the noisy bit: a biased
coin, a noisy communication channel, and a classifier's verdict. All three
share a common structure that deserves a name and a precise description.
The \emph{binary symmetric channel} is that common structure, stripped to
its essentials.

The input is a single Boolean value: $\top$ (``in the set'', ``heads'',
``positive'') or $\bot$ (``not in the set'', ``tails'', ``negative''). The
output is also Boolean. The channel transmits the input faithfully most of
the time, but occasionally flips it. The defining property of the
\emph{symmetric} channel is that the flip probability is the same
regardless of which input is being transmitted: a positive is
misreported as negative with the same probability that a negative is
misreported as positive.

This symmetry is a genuine constraint. Most real noisy Boolean systems
do not satisfy it: a spam filter is far more likely to mis-classify
legitimate mail as spam (false positive) than to miss actual spam (false
negative). The BSC is the ideal symmetric case, and understanding it
precisely is the right starting point before the asymmetric generalization
of \Cref{sec:asymmetric-channel}.

\begin{definition}[Binary Symmetric Channel]
\label{def:bsc}
A \emph{binary symmetric channel} with crossover probability $\fprate
\in (0, 1)$ is a random map $\phi : \{\top, \bot\} \to \{\top, \bot\}$
satisfying
\[
  \Pr\{\phi(x) \neq x\} \;=\; \fprate
  \quad \text{for every } x \in \{\top, \bot\}.
\]
The parameter $\fprate$ is the \emph{crossover probability}: the
probability that the output differs from the input.
\end{definition}

Two things are worth noting immediately. First, the crossover probability
is the same for both inputs: $\Pr\{\phi(\top) \neq \top\} =
\Pr\{\phi(\bot) \neq \bot\} = \fprate$. This is the symmetry. Second,
$\fprate$ plays the role of both the false-positive rate and the
false-negative rate simultaneously: in the Bernoulli framework's
notation from \Cref{sec:three-faces}, the BSC has $\fprate = \fnrate$
(using the convention that $\fprate$ is the rate of misreporting a
negative as positive and $\fnrate$ is the rate of misreporting a positive
as negative). The BSC is the special case where these two rates collapse
to a single parameter.

\subsection*{The channel matrix}

A single parameter and a prose definition are useful, but a matrix
representation is more useful still. It displays all four input-output
probabilities at once, supports direct computation, and generalizes
cleanly to the asymmetric case in \Cref{sec:asymmetric-channel} and to
multi-query systems in \Cref{ch:axioms-kronecker}.

We organize the channel as a $2 \times 2$ matrix $Q$, called the
\emph{channel transition matrix} or simply the \emph{channel matrix}. The
convention throughout this book is:

\begin{notation}
The channel matrix $Q$ for a binary channel is a $2 \times 2$ stochastic
matrix. \textbf{Rows} are indexed by the input block: row 1 corresponds to
block $B_1 = \{\top\}$ (positives), row 2 to block $B_2 = \{\bot\}$
(negatives). \textbf{Columns} are indexed by the output label: column 1
corresponds to output $\top$ (reported positive), column 2 to output
$\bot$ (reported negative). The $(i,j)$-entry of $Q$ is the probability
of observing output $j$ given that the input belongs to block $B_i$. Each
row sums to 1: $Q$ is \emph{row-stochastic}.
\end{notation}

Under this convention, the four entries of $Q$ have specific names in the
confusion-table vocabulary. For the order-2 Boolean model with $B_1 =
\{\top\}$ and $B_2 = \{\bot\}$:

\begin{center}
\begin{tabular}{l|cc}
          & output $\top$ & output $\bot$ \\
\hline
input $\top$ ($B_1$) & $\tprate$ (true-positive rate) & $\fnrate$ (false-negative rate) \\
input $\bot$ ($B_2$) & $\fprate$ (false-positive rate) & $\tnrate$ (true-negative rate) \\
\end{tabular}
\end{center}

So the general order-2 channel matrix is
\[
  Q \;=\;
  \begin{pmatrix}
    \tprate & \fnrate \\
    \fprate & \tnrate
  \end{pmatrix}
  \;=\;
  \begin{pmatrix}
    1 - \fnrate & \fnrate \\
    \fprate & 1 - \fprate
  \end{pmatrix},
\]
where we use $\tprate = 1 - \fnrate$ and $\tnrate = 1 - \fprate$ from the
row-stochastic constraint.

For the BSC, the symmetry constraint $\fprate = \fnrate$ collapses this
to a single parameter:
\[
  Q_{\mathrm{BSC}} \;=\;
  \begin{pmatrix}
    1 - \fprate & \fprate \\
    \fprate & 1 - \fprate
  \end{pmatrix}.
\]
Both rows sum to 1. Both diagonal entries equal $1 - \fprate$ (correct
output). Both off-diagonal entries equal $\fprate$ (flip). The two rows
are simply permutations of each other, which is the matrix expression of
symmetry. The BSC has a single free parameter; the general order-2 channel
has two. We return to this reduction in \Cref{sec:model-order}.

\subsection*{Reading the matrix}

Before computing with $Q$, it is worth reading it carefully for a small
$\fprate$. Take $\fprate = 0.05$ as a concrete case.

\[
  Q \;=\;
  \begin{pmatrix}
    0.95 & 0.05 \\
    0.05 & 0.95
  \end{pmatrix}.
\]

Reading row 1: a positive input ($\top$) is correctly reported as positive
($\top$) with probability $0.95$ and incorrectly reported as negative
($\bot$) with probability $0.05$. Reading row 2: a negative input ($\bot$)
is incorrectly reported as positive ($\top$) with probability $0.05$
(a false positive) and correctly reported as negative ($\bot$) with
probability $0.95$. The channel is ``mostly right'' with a small error
floor.

The channel matrix carries more information than the crossover
probability alone: it tells you the probability of every input-output
pair simultaneously, without requiring you to remember which rate is
which or convert between TPR and FNR in your head. Once we introduce
multiple queries in \Cref{sec:channel-composition} and \Cref{sec:kronecker},
the matrix form pays for itself many times over.

\subsection*{Worked example: BSC with $\fprate = 0.05$}

We verify all four probabilities and three derived quantities for
$\fprate = 0.05$.

\medskip
\noindent\textbf{The matrix.}
\[
  Q \;=\;
  \begin{pmatrix}
    1 - \fprate & \fprate \\
    \fprate & 1 - \fprate
  \end{pmatrix}
  \;=\;
  \begin{pmatrix}
    0.95 & 0.05 \\
    0.05 & 0.95
  \end{pmatrix}.
\]

\medskip
\noindent\textbf{Row sums.}
Row 1: $0.95 + 0.05 = 1.00$. Row 2: $0.05 + 0.95 = 1.00$. Both rows
sum to 1, confirming that $Q$ is row-stochastic.

\medskip
\noindent\textbf{Single-query error probability.}
If the input is drawn uniformly from $\{\top, \bot\}$, each with
probability $1/2$, the probability of an incorrect output is
\[
  \Pr\{\text{error}\}
  \;=\; \tfrac{1}{2} \cdot \Pr\{\phi(\top) \neq \top\}
       + \tfrac{1}{2} \cdot \Pr\{\phi(\bot) \neq \bot\}
  \;=\; \tfrac{1}{2} \cdot 0.05 + \tfrac{1}{2} \cdot 0.05
  \;=\; 0.05.
\]
Under a uniform prior on the input, the error probability equals the
crossover probability $\fprate$. This is not a coincidence: symmetry
means the marginal error rate does not depend on the input distribution,
as long as $\fprate = \fnrate$.

\medskip
\noindent\textbf{Two specific joint probabilities.}
Suppose the input is positive ($\top$). The probability of a correct
positive output is $Q_{11} = 0.95$. The probability of a false-negative
output is $Q_{12} = 0.05$.

Now suppose the input is negative ($\bot$). The probability of a false
positive is $Q_{21} = 0.05$. The probability of a correct negative is
$Q_{22} = 0.95$.

\medskip
\noindent\textbf{Accuracy.}
For a uniform prior, accuracy $= \frac{1}{2} \cdot 0.95 + \frac{1}{2}
\cdot 0.95 = 0.95 = 1 - \fprate$. As expected, the BSC with crossover
$0.05$ has $95\%$ accuracy under a uniform input distribution.

\subsection*{The BSC as the symmetric special case of order 2}

The source papers in the Bernoulli framework use two phrases for the BSC
that are worth aligning before proceeding. Some sections call the BSC a
``first-order'' model, while \Cref{ch:axioms-kronecker} calls it the
\emph{symmetric special case of the order-2 model}. Both are correct, but
they refer to different things: the BSC's model order, and the family we
embed it in.

Model order, formalized in \Cref{sec:model-order}, counts the blocks in
the partition of the input domain that groups together exactly those
inputs sharing a common error rate. For the BSC both inputs share the
single crossover rate $\fprate$, since $\Pr\{\phi(\top) \neq \top\} =
\Pr\{\phi(\bot) \neq \bot\} = \fprate$; the two inputs therefore fall in
one block, and the BSC has \emph{model order 1}. The general asymmetric
channel of \Cref{sec:asymmetric-channel}, with $\fprate \neq \fnrate$,
keeps the two inputs in separate blocks and has model order 2.

Calling the BSC ``the symmetric special case of the order-2 model'' refers
not to its model order but to the family we work in. We can always write
the BSC with the two-block discrete partition $B_1 = \{\top\}$, $B_2 =
\{\bot\}$ and then impose the constraint $\fprate = \fnrate$; chapters
2--4 work in this two-block frame throughout, because it is the natural
one for the Boolean domain. So the BSC is genuinely model order 1, studied
as a constrained member of the order-2 family.

Model order and parameter count are nonetheless distinct in general, and
\Cref{sec:model-order} separates them carefully: the Z-channel, with
$\fnrate = 0 \neq \fprate$, has model order 2 (two blocks at different
rates) yet only one free parameter. For the BSC the two counts happen to
coincide at 1; for the general asymmetric channel they coincide at 2.

\Cref{sec:axiom-2} makes the connection to partition blocks explicit in
the context of Axiom 2. The convention $B_1 = \{\top\}$, $B_2 =
\{\bot\}$ used there matches the row labeling of $Q$ throughout this
chapter and in \Cref{ch:axioms-kronecker}.

\subsection*{Connection to chapter 1: the three faces revisited}

\Cref{sec:three-faces} presented three framings of the same abstract
object. We can now describe each one precisely using the BSC channel
matrix.

\medskip
\noindent\textbf{The biased coin.} A coin with heads probability
$1 - \fprate$ and tails probability $\fprate$, where ``heads'' is the
correct output, is a BSC. The two sides of the coin are the two elements
of $\{\top, \bot\}$; the flip probability is $\fprate$ for both sides,
because the coin is symmetric. The channel matrix is exactly
$Q_{\mathrm{BSC}}$ above. We generalize to an asymmetric coin, with
different head-bias and tail-bias, in \Cref{sec:asymmetric-channel}.

\medskip
\noindent\textbf{The noisy communication channel.} This is Shannon's BSC
in its original form \cite{shannon-1948}: a bit sent over a channel where
each transmitted bit is flipped independently with probability $\fprate$.
The input alphabet is $\{0, 1\}$ and the channel matrix is $Q_{\mathrm{BSC}}$
with the mapping $1 \leftrightarrow \top$ and $0 \leftrightarrow \bot$.

\medskip
\noindent\textbf{The classifier verdict.} A binary classifier with equal
false-positive and false-negative rates is a BSC. In practice, classifiers
rarely have $\fprate = \fnrate$; the symmetric case is a simplifying
assumption. When the two rates are unequal, the classifier is better
modeled as an asymmetric binary channel, the subject of
\Cref{sec:asymmetric-channel}.

The unifying point: all three framings share the same channel matrix
$Q_{\mathrm{BSC}}$ with the same single crossover parameter $\fprate$.
The names differ; the probability model is identical.

\subsection*{What makes symmetry a strong constraint}

The BSC's symmetry constraint, $\fprate = \fnrate$, is restrictive but
often useful as a first approximation. Three properties follow from it.

First, the channel matrix $Q_{\mathrm{BSC}}$ is a \emph{circulant}
matrix: its rows are cyclic permutations of each other. This makes
the algebra of composing BSCs, discussed in \Cref{sec:channel-composition},
particularly clean.

Second, the \emph{capacity} of the BSC, defined as the maximum
mutual information between input and output, is $1 - H_2(\fprate)$ bits
per use, where $H_2(\fprate) = -\fprate \log_2 \fprate - (1-\fprate)
\log_2(1-\fprate)$ is the binary entropy function. Capacity is maximized
at $\fprate = 0$ (a perfect channel, transmitting one bit per use without
error) and at $\fprate = 1$ (a perfectly anti-correlated channel, which
also transmits one bit per use after relabeling outputs); it is minimized
at $\fprate = 1/2$, where the output is statistically independent of the
input and no information is transmitted. Exercise~\ref{ex:bsc-capacity}
works through the derivation.

Third, the BSC is its own \emph{inverse channel} in the following sense:
composing a BSC with crossover $\fprate$ with another BSC with the same
crossover $\fprate$ yields a BSC with crossover $2\fprate(1-\fprate)$,
which is strictly greater than $\fprate$ for $0 < \fprate < 1/2$. Two
imperfect copies of the same symmetric channel make things worse, not
better. The formula for composed BSC crossovers is derived in full in
\Cref{sec:channel-composition}.

\subsection*{Notation in this chapter}

Throughout chapters 2 and beyond, the following convention applies:
\begin{itemize}
  \item $Q$ always denotes a channel transition matrix, with rows indexed
    by input block and columns indexed by output label.
  \item $\fprate$ is the false-positive rate: the probability that a negative
    input is reported as positive.
  \item $\fnrate$ is the false-negative rate: the probability that a positive
    input is reported as negative.
  \item $\tprate = 1 - \fnrate$ is the true-positive rate, and
    $\tnrate = 1 - \fprate$ is the true-negative rate.
  \item For the BSC, $\fprate = \fnrate$, and we use $\fprate$ for the
    single crossover parameter throughout.
\end{itemize}

The source papers in this research program sometimes use (FPR, TPR) rather
than (FPR, FNR) as the parameter pair. The two choices are equivalent:
TPR $= 1 - \text{FNR}$ is a deterministic function, so any statement about
one pair translates directly to the other.\footnote{Specifically,
the companion paper \cite{towell-bernoulli-sets} uses (FPR, TPR) in places.
In converting source material to this textbook, we have replaced that
convention with $(\fprate, \fnrate)$ throughout.}

\subsection*{Looking ahead}

The BSC is the symmetric special case of a more general object. In
\Cref{sec:asymmetric-channel} we lift the constraint $\fprate = \fnrate$
and allow the two error rates to be independent: the resulting
\emph{asymmetric binary channel} covers Bloom filters (where $\fnrate = 0$
by construction) and medical tests (where FPR and FNR can be tuned
independently by adjusting a threshold). In \Cref{sec:channel-composition}
we ask what happens when two channels are applied in sequence, and derive
the composed channel's matrix from the two input matrices via matrix
multiplication. In \Cref{sec:model-order} we formalize what we have been
calling ``order 2'' and prove the Boolean ceiling theorem: because the
Boolean domain has exactly two elements, every binary channel has model
order at most 2.
